% ch9_a §9.1–9.2 (书页 292–300 / p306–p314)
\chapter{费米面}

\epigraphbook{你以这样一副可疑的形貌前来。}{Hamlet}

\section{强磁场}

一般说来，高温下金属的电子性质是对整个费米面上的速度、波函数、跃迁概率等量所作的复杂平均。除了由 Kohn 效应（§5.4）和正电子湮没（§5.8）得到的一些模糊而零碎的信息之外，要根据实验观测真正描绘出这个面，就得使用纯净的样品，并在低温下工作，使电子的散射不至于把现象弄模糊。当电子的弛豫时间非常短时，费米面是否还能恰当地说成“存在”，甚至都是可疑的。例如，平均自由程为 $\Lambda$ 的电子，其动量的不确定度约为 $\hbar/\Lambda$ 的量级；在高温下，这个不确定度很可能大到足以同倒格子空间中能量面的某些细致结构相比拟。

我们已经讨论过反常趋肤效应（§8.7）。关于费米面的几乎所有其他研究都有赖于强磁场的使用。磁场对一个电子态的作用由式 (6.40) 给出：
\begin{equation*}
\dot{\mathbf{k}}=\frac{e}{c\hbar}\,\mathbf{v}\wedge\mathbf{H}. \tag{9.1}
\end{equation*}
这意味着矢量 $\mathbf{k}$ 的变化是
\begin{itemize}
\item[(i)] 垂直于 $\mathbf{H}$ 的方向；
\item[(ii)] 垂直于 $\mathbf{v}$，而 $\mathbf{v}$ 本身又垂直于能量面。
\end{itemize}
因此，$\mathbf{k}$ 必然被限制在\emph{轨道}（orbit）上，即由费米面同垂直于 $\mathbf{H}$ 的平面相交所得的曲线。磁场只是驱使代表点沿这条轨道环行，而不改变能量。

如果电子不被散射，它环行一周所经历的时间为
\begin{equation*}
\frac{2\pi}{\omega_H}=\frac{c\hbar}{eH}\oint\frac{dk}{v_{\perp}}, \tag{9.2}
\end{equation*}
其中 $v_{\perp}$ 是 $\mathbf{v}$ 在垂直于 $\mathbf{H}$ 的平面内的分量，取在 $\mathbf{k}$ 点处。
相应的频率 $\omega_H$ 称为回旋频率（cyclotron frequency）。对于自由电子，由初等几何可得
\begin{equation*}
\oint\frac{dk}{v_{\perp}}=\frac{m}{\hbar}\oint\frac{dk}{k_{\perp}}=\frac{2\pi m}{\hbar}, \tag{9.3}
\end{equation*}
于是
\begin{equation*}
\omega_H=\frac{eH}{mc}. \tag{9.4}
\end{equation*}
%FIG{153}: {磁场中电子的“轨道”。}
%FIG{154}: {}

习惯上还要再定义另一种“有效质量”——回旋质量（cyclotron mass）$m_H^{*}$，使得
\begin{equation*}
\omega_H=\frac{eH}{m_H^{*}c}. \tag{9.5}
\end{equation*}
应当注意，这个质量并不就是电子的动力学质量。它是轨道的属性，而不是某个特定电子态的属性。

回旋质量有一个有用的几何定义，它来自式 (6.2) 中隐含的一个事实：
\begin{equation*}
v_{\perp}=\frac{1}{\hbar}\frac{d\mathcal{E}}{dk_{\perp}}, \tag{9.6}
\end{equation*}
其中 $dk_{\perp}$ 是 $k$ 在轨道平面内、沿垂直于费米面方向的增量。由式 (9.2) 和式 (9.5) 可得（图 154）
\begin{align*}
m_H^{*}&=\frac{eH}{c}\,\frac{1}{2\pi}\,\frac{c\hbar}{eH}\oint\frac{dk}{v_{\perp}}\\
&=\frac{\hbar^{2}}{2\pi}\oint\frac{dk_{\perp}}{d\mathcal{E}}\,dk\\
&=\frac{\hbar^{2}}{2\pi}\frac{\partial\mathcal{A}}{\partial\mathcal{E}}, \tag{9.7}
\end{align*}
其中 $\mathcal{A}$ 是轨道在垂直于 $\mathbf{H}$ 的平面内所围的面积。

\section{回旋共振}

自然会想到：电子的这种周期运动可以通过与频率合适的电磁场发生共振而被探测到。唯一的条件是，电子在被散射之前至少要沿轨道环行一周。要为系统对沿磁场轴向传播的圆偏振电磁波的响应建立一个形式理论，是相当容易的。

例如，我们用一组新变量来标记 k 空间。一个点 $\mathbf{k}$ 按如下方式确定：
\begin{itemize}
\item[(i)] 能量 $\mathcal{E}(\mathbf{k})$；在磁场中它保持不变。
\item[(ii)] $\mathbf{k}$ 沿磁场方向的分量 $k_H$；它也保持不变。
\item[(iii)] 一个“相变量” $\phi$，由沿轨道的、与式 (9.2) 中出现的同一种积分来定义，即
\end{itemize}
\begin{equation*}
\phi=\omega_H\frac{c\hbar}{eH}\int^{\mathbf{k}}\frac{dk}{v_{\perp}}. \tag{9.8}
\end{equation*}
这是一个便利的变量，因为在磁场作用下它自动以恒定速率增大：
\begin{equation*}
\dot{\phi}=\omega_H, \tag{9.9}
\end{equation*}
且当环行一整圈时 $\phi=2\pi$。

我们可以把这些变量用作 Boltzmann 方程的坐标框架。可以把式 (7.10) 中的“失衡”部分 $g_{\mathbf{k}}$ 表示为 $(\mathcal{E},k_H,\phi)$ 的函数。如果像式 (8.84) 中那样假定一个弛豫时间，那么只需再添一项
\begin{equation*}
\left.\frac{\partial f}{\partial t}\right]_{\text{mag.}}=\dot{\phi}\,\frac{\partial g}{\partial\phi}=\omega_H\frac{\partial g}{\partial\phi} \tag{9.10}
\end{equation*}
就能把磁场对电子分布的影响包括进来。这与式 (7.137) 中所用的表达式等价，只是更为紧凑，在那里磁场被假定是小的。

我们取一个频率为 $\omega$ 的随时间变化的场，并忽略 $\mathbf{E}$ 随空间的变化。再仿照式 (8.86)，假定
\begin{equation*}
g(\mathcal{E},k_H,\phi)=\left(-\frac{\partial f^0}{\partial\mathcal{E}}\right)\Phi(k_H)\,e^{i(\phi-\omega t)}, \tag{9.11}
\end{equation*}
即电子分布（即便不是每一个单个电子）实际上以外场的频率被驱动沿轨道环行。于是我们的 Boltzmann 方程
\begin{equation*}
e\mathbf{E}\cdot\mathbf{v}\left(-\frac{\partial f^0}{\partial\mathcal{E}}\right)=\frac{g}{\tau}+\omega_H\frac{\partial g}{\partial\phi}+\frac{\partial g}{\partial t} \tag{9.12}
\end{equation*}
具有解
\begin{equation*}
\Phi(k_H)=\frac{e\tau\,\mathbf{v}\cdot\mathbf{E}}{1+i(\omega_H-\omega)\tau}, \tag{9.13}
\end{equation*}
和先前一样，偏离正比于场的强度，但除非外加频率与系统的固有回旋频率相同，它将带着相位差沿轨道跟随转动。

正如式 (8.90) 那样，金属的行为就好像它具有如下电导率：
\begin{equation*}
\sigma(\omega)=\sigma(0)\,\frac{1-i(\omega_H-\omega)\tau}{1+(\omega_H-\omega)^{2}\tau^{2}}. \tag{9.14}
\end{equation*}
观测材料的表面电阻、反射本领、吸收等，就会在频率 $\omega=\omega_H$ 处看到一条宽度为 $1/\tau$ 的回旋共振线。这一结果与式 (8.51) 中的光学共振理论之间的类似是显然的。对自由电子，通过对运动方程的简单分析不难得到这个结果——但我们在这里证明的结果要略微更普遍一些，尽管对于任意费米面它并不完全精确，因为在那里的假设 (9.11) 忽略了变量 $\phi$ 的高次谐波。

这一公式是对旋转方向与磁场中电子固有回旋运动\emph{相同}的圆偏振波束而言的。
显而易见，从几何上的考虑即可知道：偏振方向相反的波束，其公式将以 $+\omega$ 代替 $-\omega$；这样就不会发生共振。但考虑 $\sigma(\omega)$ 的\emph{虚部}，它对复折射率 (8.7) 的\emph{实部}有贡献，从而改变电磁波的相速。对于沿磁场方向传播的线偏振波束，两个圆偏振分量以不同的速度行进，因此当波穿过晶体时，偏振面便不断旋转。这就是众所周知的 Faraday 效应。

对半导体而言，这几乎就是描述回旋共振所需的全部理论。载流子无论电子还是空穴，都占据着 $\mathcal{E}(\mathbf{k})$ 函数的一些小口袋，它们位于能带的极小或极大处。在这些邻域内，能量面通常是椭球面（参见 §2.5）；可以证明，$\omega_H$ 在椭球的所有平行截面上都是相同的，因此能观测到式 (9.14) 所预言的谱线。当然，$\omega_H$ 随取向而变化，由此可以提供关于椭球各轴的信息。在区中心价带极大附近情形更为复杂，那里能级被自旋–轨道相互作用（§3.9）所分裂，但分析基本上仍循上述路线进行。

对金属来说，有两个似乎横亘在前的困难。为了分辨共振线，我们需要在高频下工作，使 $\omega\tau\gg 1$。在这些条件下，趋肤深度非常小——实际上我们已处在 §8.7 的反常区域。于是，磁场中电子的螺旋轨迹在真实空间中的半径，将远大于电场所能穿透的距离。

这提示我们应当把磁场垂直于金属表面施加，这样发生共振的电子就是趋肤层内的“有效”电子，能够吸收功率。但对这种情形作一分析就会发现，“共振”被抹平到几乎无法探测。把式 (9.14) 代入式 (8.7)，原则上可以看出这是如何发生的。正如式 (8.91) 那样，复折射率有很大的虚部，主要原因在于 $N^{2}$ 的实部很大且为负。金属的电导太高了；我们所能测量的只是其反射本领的微小亏损，如式 (8.94) 那样。要看到在半导体中很容易探测到的那种共振线，我们就必须工作在等离体频率 $\omega_p$ 之上，而这将需要极其巨大的磁场。

另一方面，假设我们手头有“好的”回旋材料，例如 4°K 的极纯 Na，在几千奥斯特的磁场中其 $\omega_H\tau\gg 1$。如果我们以非常低的频率工作——每秒不过几个周期——就可以使折射率 (8.7) 接近于实数，并观察电磁波穿过一块金属板的传播。这可由式 (9.14) 得出；令 $\omega\ll\omega_H$，我们得到
\begin{align*}
\mathsf{N}&=\left(\epsilon+\frac{4\pi i\sigma(\omega)}{\omega}\right)^{1/2}\\
&\approx\left(\frac{4\pi\sigma(0)}{\omega\omega_H\tau}\right)^{1/2}=\left(\frac{4\pi nec}{H\omega}\right)^{1/2}=\frac{\omega_p}{(\omega\omega_H)^{1/2}}. \tag{9.15}
\end{align*}
结果表明这是一个非常大的数，约为 $10^{10}$ 的量级。因此波在金属中仅以每秒几厘米的速度行进——一种以“人”的长度和时间尺度来衡量的现象。在式 (9.15) 中，重要的是 $\omega$ 与 $\omega_H$ 应当同号；式 (9.8) 和式 (9.11) 确立的约定意味着，我们所处理的是与电子回旋运动同相位旋转的圆偏振波。这样，一个螺旋波（helicon）模就把电子气的回旋激发与等离激元激发结合在一起。遗憾的是，从这种现象所能获得的信息是有限的，因为式 (9.15) 在一级近似上并不仅仅依赖于费米面的几何。在补偿半导体（§4.6）中，电子与空穴的螺旋波项彼此抵消，只剩下描述 Alfvén 波传播的高阶项——这是天体物理中另一种人们熟悉的磁–等离体现象。

事实证明，当磁场\emph{平行}于表面时可以探测到很强的共振。这就是 Azbel'–Kaner 共振，其成因如下。电子沿螺旋轨道回旋，其轴向（比如说沿 $x$ 方向）与表面平行。每回旋一周，它都进入趋肤深度一次，“看见”振荡的电场。如果它每次返回时都与场同相位，它就会从场中获得能量，从而观察到“共振”。

推导一个描述这种现象的公式并不困难。我们把普遍的动理学表达式 (8.101) 写成如下形式：
\begin{equation*}
\mathbf{J}(\mathbf{r},t)=\frac{e^{2}}{4\pi^{3}\hbar}\int\mathbf{v}\left\{\int_{-\infty}^{t}\mathbf{v}(\mathbf{r}',t')\cdot\mathbf{E}(\mathbf{r}',t')\,e^{-(t-t')/\tau}\,dt'\right\}\frac{dS_F}{v}. \tag{9.16}
\end{equation*}
这里我们已经对费米面积分，得到总电流。但我们仍然明显地保留了论证的主要特征
%FIG{155}: {电场使费米分布发生位移。这个位移被驱动沿费米面环行，并随时间衰减。}
%FIG{156}: {Azbel'–Kaner 共振。}
——即速度为 $\mathbf{v}$ 的电子在时刻 $t$ 对 $\mathbf{r}$ 处电流的贡献，取决于场在早先时刻 $t'$ 于电子曾经过的点 $\mathbf{r}'$ 处所给予它的冲量；不过对这个冲量的记忆随时间消退，由一个以弛豫时间 $\tau$ 为特征的衰减因子支配。这种对电传导的描述，在精神上与 Kubo 公式 (7.15) 非常接近。

现在的轨迹是螺旋形的。电子沿一条弯曲路径到达 $\mathbf{r}$，沿这条路径，随时间的衰减仍由同一个指数因子支配。在一个周期内，该因子将为 $\exp(-2\pi/\omega_H\tau)$。电场只在电子处于趋肤深度内的短暂时刻起作用；每回旋一周，电子相对振荡场的相位就将增减一个量 $(2\pi\omega/\omega_H)$。于是，我们可以把式 (9.16) 中的时间积分化为第一次穿过趋肤深度所得的结果，再乘以随后历次穿过的衰减因子之和，即
\begin{equation*}
(1+e^{-w}+e^{-2w}+\ldots)=\frac{1}{1-e^{-w}}, \tag{9.17}
\end{equation*}
其中
\begin{equation*}
e^{-w}=\exp\left(-\frac{2\pi}{\omega_H\tau}-i\,\frac{2\pi\omega}{\omega_H}\right). \tag{9.18}
\end{equation*}
%FIG{157}: {电子在角 $\phi_m$ 内是“有效”的：(a) 在真实空间的螺旋路径上；(b) 在费米面上的轨道上。}
我们可以估计电子的“有效性”：注意到每个电子停留在趋肤深度内的时间只有 $\phi_m/\omega_H$ 的量级，于是（图 157）
\begin{align*}
\phi_m^{2}&\approx\frac{\delta'}{R}\\
&\approx\frac{\delta'\omega_H}{v_x}, \tag{9.19}
\end{align*}
其中 $R$ 是真实空间中轨道的半径，而 $v_x$ 是电子在垂直于磁场的截面平面内、平行于表面方向的速度。

再者，大多数穿过趋肤深度的电子会打到金属表面上，在那里被散射而损失掉。只有在趋肤深度内平行于表面运动的电子才是有效的。用式 (8.107) 中那类论证，我们只计入属于下述费米面元面积的那些电子：
\begin{equation*}
dS=2\beta\phi_m|\rho_y|\,dk_y, \tag{9.20}
\end{equation*}
其中 $\rho_y$ 同样是轨道的曲率半径，取在厚度为 $dk_y$ 的切片上、速度平行于金属表面的那一点处。

把这些论证合在一起，我们得到一个本质上与式 (8.109) 相像的结果，即
\begin{equation*}
\sigma'=\frac{e^{2}\beta\delta'}{2\pi^{3}\hbar}\int\frac{|\rho_y|\,dk_y}{1-e^{-w}}. \tag{9.21}
\end{equation*}
把式 (9.21) 代入式 (8.104) 和式 (8.105)，表面电阻便可算出。

就我们目前的目的而言，有意义的是共振因子 (9.17)。它是振荡的，当
\begin{equation*}
\omega=n\omega_H \tag{9.22}
\end{equation*}
时出现峰值。

实践中，人们保持频率 $\omega$ 不变，而改变磁场。于是，由式 (9.5)，
\begin{equation*}
\frac{1}{H}=n\,\frac{e}{\omega m_H^{*}c} \tag{9.23}
\end{equation*}
就是 $H$ 变化时出现峰值的条件。值得注意的是，Azbel'–Kaner 共振并不像普通回旋共振 (9.14) 那样只有一个单一的峰。只要频率使得电磁场在电子相继两次出现于表面的时刻之间恰好经历整数个振荡周期，能量就仍会被馈入。每个共振的尖锐程度显然取决于式 (9.18) 中 $w$ 的实部的大小；要有强的效应，我们需要 $\omega_H\tau\gg 1$。

还有更进一层的复杂情况。$\omega_H$ 的值（或者等价地，$m_H^{*}$ 的值）依赖于轨道。对于一般的费米面，它在不同的轨道上并\emph{不}相同，即使这些轨道属于垂直于同一磁场的一组平行切片。这样，式 (9.21) 中的变量 $w$ 本身就是 $k_y$ 的函数；来自不同切片的贡献未必保持同相位。

共振曲线形状的精确计算到此变得非常复杂。但人们可以论证——而且这一论证还可以用形式的数学加以支持——主要贡献将来自那些 $m_H^{*}$ 作为 $k_y$ 的函数取平稳值（极大或极小）的截面。如果从每个这样的点出发、把 $m_H^{*}$ 作为 $\delta k_y$ 的函数展开，常数项就提供一个共振峰，而展开式的次一项（$(\delta k_y)^{2}$ 量级）并不能把它消除。来自费米面其他区域——那里 $m_H^{*}$ 随 $\delta k_y$ 线性变化——对积分的贡献则会失去相位而相互抵消。我们遇到的正是波动光学中的惯常论证，那里 Cornu 螺线是对求解很有用的一种几何构图。

%FIG{158}: {极值轨道。}
